% ============================================================
% Chapter 4 — System Architecture and Data Model
% ============================================================
\chapter{System Architecture and Data Model}
\label{ch:architecture}

\roadmap{This chapter documents the delivered system as software: the
monorepo layout, the backend and frontend stacks, the out-of-process
computer-vision workers, the storage and job-orchestration fabric, the
authentication and role model, and the database schema of roughly ninety
models. It closes with the platform's engineering-maturity profile,
including the known weak points, so that later chapters can cite rather
than restate them.}

\section{Overview and monorepo layout}
\label{sec:arch-overview}

The platform is a TypeScript monorepo with two Python workers. Its parts
are:

\begin{itemize}
  \item \code{apps/api} --- the backend: a NestJS application organized
  into 46 feature modules over Prisma/PostgreSQL, with Redis for queues
  and rate limiting and MinIO (S3-compatible) for binary storage.
  \item \code{apps/web} --- the frontend: a React/Vite single-page
  application with role-gated routes for students, teachers, and
  administrators.
  \item \code{apps/pyfeat-worker} and \code{apps/openface3-worker} ---
  Python computer-vision workers consuming Redis job queues and writing
  results to PostgreSQL (\cref{ch:sensing}).
  \item \code{packages/shared} and \code{packages/config} --- shared
  types and defaults, including the affective-mapping rule defaults used
  by both backend and frontend.
  \item \code{gals-studio} --- a separate researcher-facing studio server
  hosting cohort-level export and coding tools referenced in
  \cref{sec:sensing-export}.
  \item \code{analysis} --- offline Python export scripts invoked by the
  API's job endpoints.
\end{itemize}

\begin{figure}[htbp]
\centering
\begin{tikzpicture}[
  node distance=7mm and 9mm,
  box/.style={draw, rounded corners=2pt, align=center, minimum height=9mm,
              inner sep=5pt, font=\small},
  store/.style={box, fill=blockgreen},
  svc/.style={box, fill=blockblue},
  wrk/.style={box, fill=blockamber},
  arr/.style={-{Stealth[length=2.4mm]}, semithick}
]
  \node[svc] (web) {React/Vite web app\\\scriptsize student, teacher, admin surfaces};
  \node[svc, below=of web] (api) {NestJS API (\code{apps/api})\\\scriptsize 46 modules, Prisma ORM};
  \node[store, below left=9mm and 2mm of api] (pg) {PostgreSQL\\\scriptsize $\sim$90 models};
  \node[store, below=9mm of api] (redis) {Redis\\\scriptsize queues, throttling};
  \node[store, below right=9mm and 2mm of api] (minio) {MinIO (S3)\\\scriptsize video, PDFs, exports};
  \node[wrk, right=14mm of redis] (of3) {OpenFace3 worker\\\scriptsize 8-class emotion};
  \node[wrk, below=5mm of of3] (pyf) {Py-Feat worker\\\scriptsize 18 AU intensities};
  \node[svc, left=12mm of pg] (studio) {Studio server\\\scriptsize cohort exports, coding};

  \draw[arr] (web) -- node[right, font=\scriptsize]{REST + WebSocket} (api);
  \draw[arr] (api) -- (pg);
  \draw[arr] (api) -- (redis);
  \draw[arr] (api) -- (minio);
  \draw[arr] (redis) -- node[above, font=\scriptsize]{job queues} (of3);
  \draw[arr] (redis.east) to[out=-15,in=180] (pyf.west);
  \draw[arr] (of3.south west) to[out=200,in=-20] node[below, font=\scriptsize]{results} (pg.south east);
  \draw[arr] (pyf.west) to[out=190,in=-40] (pg.south);
  \draw[arr] (minio.south) to[out=-60,in=-60, looseness=1.4] node[below, font=\scriptsize]{video segments} (pyf.south east);
  \draw[arr] (studio) -- (pg);
\end{tikzpicture}
\caption[Full-stack architecture]{Full-stack architecture. The web client
speaks REST and WebSocket to the NestJS API; PostgreSQL holds all
relational and JSON data; MinIO holds binaries (webcam segments, PDFs,
export bundles); Redis carries the computer-vision job queues consumed by
the two Python workers, which write frame-level results back to
PostgreSQL. The studio server reads the same database for cohort-level
research exports.}
\label{fig:architecture}
\end{figure}

\section{Backend}
\label{sec:arch-backend}

The API's 46 modules fall into five functional groups, inventoried
exhaustively in the project's evidence
pack:\footnote{\code{docs/capstone-evidence/01_functionality_inventory.md}}
instructor-side authoring and curation (courses, modules, questions,
assessments, generation, knowledge components, evaluation, coverage,
publish gate, RAG, user management); student-side learning (attempts,
interventions, dialogue, notes, chat history, mastery, student RAG);
sensing and logging ingest (logs, recording, webgazer, pupil-size,
activity-log); analysis and replay (analytics, affective-mapping,
text-mining, replay-annotations); and infrastructure (auth, blob,
event-bus, jobs, health, LLM registry, pre-generation, and the CV-job
orchestrators). Cross-cutting conventions include a global exception
filter normalizing Zod and Prisma errors, Zod validation pipes on the
newer modules, Redis-backed rate limiting, and per-call LLM cost
accounting into \dbtable{LlmUsageLog} against a
\dbtable{LlmModelPricing} table.

Asynchronous work uses two mechanisms. A PostgreSQL transactional outbox
(\dbtable{EventQueue}, claimed with \code{FOR UPDATE SKIP LOCKED}) carries
domain events such as grading requests; in-process pollers fan results
out over WebSocket. Heavier computer-vision work is queued to Redis lists
consumed by the Python workers. Two engineering caveats are recorded here
once and assumed thereafter: the outbox has no retry or dead-letter
mechanism (failed events are marked and left), and several in-process
schedulers (pre-generation, pollers) assume a single API
instance.\footnote{\code{apps/api/src/event-bus/event-bus.service.ts};
\code{apps/api/src/pre-generation}}

\section{Frontend}
\label{sec:arch-frontend}

The web client organizes routes under role guards: teacher/admin surfaces
(course builder, studio, question tools, review hub) and student surfaces
(dashboard, course views, dialogue workspace, assessments, review queue),
with all student learning routes wrapped in a biometrics provider that
gates rendering on capture permissions. Two component families matter for
later chapters: the \emph{FloatingChatbot} family (floating and docked
variants sharing one chat panel and the four intervention views), and the
\emph{student-logs} family (the session list, the tabbed session viewer,
and the replay engine with its AOI scoring and export libraries). The
teacher landing dashboard is a placeholder (\cref{sec:scope-item8});
every other routed surface is functional.

\section{The computer-vision workers}
\label{sec:arch-workers}

Facial analysis is deliberately out-of-process. When a webcam segment
upload completes, the recording service enqueues one job per enabled
pipeline; each worker \code{blpop}s its Redis list, fetches the segment
from MinIO, samples frames at its configured rate, runs genuine model
inference, and bulk-inserts frame rows. The OpenFace~3.0 worker loads the
RetinaFace alignment and multitask weights at boot and fails fast if they
are absent; the Py-Feat worker instantiates the library's detector with
the \code{xgb} AU model. Neither returns placeholder data; both are inert
unless deployed with their weights, which is an operational rather than
scientific caveat. \Cref{ch:sensing} specifies exactly which worker
produces which signal, since this attribution was historically confused
in project documentation.

\section{Storage and job orchestration}
\label{sec:arch-storage}

PostgreSQL is the single source of truth for relational data, JSON
payloads (intervention session state, snapshot AOIs, rule sets), and all
frame-level sensing results. MinIO holds webcam segments, uploaded PDFs,
page images for multimodal retrieval, and generated export bundles, with
presigned URLs brokered by the API. Redis carries the CV job queues and
the rate limiter. Embeddings are stored as JSON arrays on chunk rows and
compared in memory---adequate at study scale, and flagged as a
scalability limit in \cref{ch:limitations}.

\section{Authentication and roles}
\label{sec:arch-auth}

Users hold one of three roles (student, teacher, admin). Authentication
accepts email or a teacher-assigned login identifier and issues a 24-hour
JWT; registration and login are rate-limited; teachers provision students
in bulk with temporary passwords. Authorization is enforced by role
guards at the route level and ownership checks in services (e.g.,
researchers may only read their own annotation codes). Consent for webcam
recording is stored per student per course and checked before capture
(\cref{ch:sensing}).

\section{Data model}
\label{sec:arch-datamodel}

The Prisma schema defines approximately ninety models in eight domains;
\cref{tab:arch-domains} summarizes them, and
\cref{app:schema} reproduces the core research tables verbatim. Three
schema-level facts recur in later chapters. First, all sensing streams
anchor on \dbtable{StudentSession}, and one
\dbtable{session_sync_anchors} row per session ties wall-clock,
monotonic, and server-receive time together---the platform's single time
base. Second, deletion semantics are deliberate: purging a session
cascades to its logs, snapshots, and frames, but chatbot conversation
rows survive with a nulled session key, and researcher annotations
cascade with the session (\cref{sec:sensing-coding}). Third, two fields
document intent rather than behavior---the BKT parameter columns on
\dbtable{UserMastery} and the head-pose columns on
\dbtable{EmotionFrame} are present but never populated---and this report
cites them only as such.

\begin{table}[htbp]
\centering
\small
\caption{Data-model domains (Prisma schema, $\sim$90 models).}
\label{tab:arch-domains}
\begin{tabular}{@{}p{0.3\linewidth}p{0.62\linewidth}@{}}
\toprule
Domain & Representative models \\
\midrule
Identity \& courses & \dbtable{User}, \dbtable{Course},
\dbtable{Enrollment}, \dbtable{CourseModule}, \dbtable{ModuleItem},
\dbtable{PageContentVersion} \\
Assessment & \dbtable{Question}, \dbtable{Assessment},
\dbtable{Attempt}, \dbtable{GradingResult}, \dbtable{KcEvidence} \\
RAG corpora & \dbtable{SourceDocument}, \dbtable{DocumentChunk},
\dbtable{StudentSourceDocument}, \dbtable{StudentRagChunk} \\
Dialogue \& chatbot & \dbtable{DialogueSession},
\dbtable{DialogueMessage}, \dbtable{ChatbotMessage},
\dbtable{StudioOutput}, \dbtable{DialogueNote} \\
Knowledge components & \dbtable{ProposedKC}, \dbtable{KcEdge},
\dbtable{KcContentMapping}, \dbtable{KnowledgeVersion};
\dbtable{KnowledgeComponent}, \dbtable{UserMastery} \\
Interventions & \dbtable{LearningIntervention},
\dbtable{SpacedRepetitionCard}, \dbtable{InterventionPromptConfig},
\dbtable{SavedInterventionReview}, \dbtable{PreGeneratedExercise} \\
Sensing \& logging & \dbtable{StudentSession}, \dbtable{ActivityLog},
nine raw interaction tables, \dbtable{SessionReplaySnapshot},
\dbtable{session_sync_anchors}, \dbtable{RecordingSegment},
\dbtable{WebgazerLog}, \dbtable{PupilSizeLog},
\dbtable{PyfeatAuResult}, \dbtable{EmotionFrame} \\
Analysis \& coding & \dbtable{AffectiveMappingConfig},
\dbtable{AffectiveStateWindow}, \code{derived\_*} tables,
\dbtable{aligned_frames}, \dbtable{EfDetection},
\dbtable{EfConstructPrompt}, \dbtable{ReplayCode},
\dbtable{ReplayAnnotation} \\
\bottomrule
\end{tabular}
\end{table}

\section{Engineering-maturity profile}
\label{sec:arch-maturity}

The evidence-pack audit rated every module and surface. The headline
profile: the intervention service, sensing ingest, replay engine, and
authoring suite are production-grade; exactly one endpoint is a declared
stub (text-mining session
reprocessing\footnote{\code{apps/api/src/text-mining/text-mining.controller.ts:255}});
the AI-grading path for structured questions writes a placeholder score
awaiting an external grader; the teacher dashboard is a placeholder; and
one security flaw---an unvalidated session identifier interpolated into a
shell command in the export-jobs
controller\footnote{\code{apps/api/src/jobs/jobs.controller.ts:18}}---is
acceptable only because the endpoint is role-restricted, and is listed
for remediation in \cref{ch:future}. These facts are carried into
\cref{ch:limitations} rather than repeated per chapter.

\medskip
\noindent With the machinery established, the next two chapters present
what learners actually experience: the two learning modes and the four
interventions.
